% Quarantine-only candidate proof.
% Date: 29 August 2026.
% This file is not a live-book input and makes no novelty or priority claim.

\begin{proof}[Proof of
\cref{lem:p213-arbitrary-valuation-source-localization}]
Use the ordinary local moment description in
\cite[Proposition~15.1.1(b)]{CLD25}.  After shrinking about $y$ inside
$W$,
there are an analytic moment $F_{\mathrm{form}}$ and a smooth superform
$\eta$ such that $\alpha=F_{\mathrm{form}}^*\eta$.  Shrink once more so
that the prescribed finite block $G$ is defined.  When $\chi$ is present,
\cite[Corollary~15.1.2]{CLD25} gives, after one further shrinking, an
analytic moment $H$ and a smooth function $\Phi$ such that
$\chi=\Phi\circ H_{\mathrm{trop}}$.  Denote this common open neighborhood
by $O\subset W$.

Apply \cref{lem:p213-normal-curve-nonrigid-localization} to $Y$ at $y$,
inside $O$.  It gives a compact analytic neighborhood $V$ of $y$ and
a moment $F_{\mathrm{bd}}$ such that
\[
 F_{\mathrm{bd,trop}}(y)\notin
 F_{\mathrm{bd,trop}}(\partial(V/O)).
\]
The inclusion $V\subset O\subset W$ is a composite of analytic-domain
inclusions.  Since $O$ is open in $W$, the relative interiors of $V$ in
$O$ and in $W$ agree; hence
$\partial(V/O)=\partial(V/W)$.  We use this common relative boundary
below.
Concatenate $F_{\mathrm{bd}}$, $F_{\mathrm{form}}$, $G$, and $H$.
Projection to the first block proves
\eqref{eq:p213-arbitrary-source-local-boundary}; projection to the form
and cutoff blocks proves
\eqref{eq:p213-arbitrary-source-local-form} and
\eqref{eq:p213-arbitrary-source-cutoff-factorization}.  Equality of two
full tropical values implies equality of their $H$-projections, so
applying $\Phi$ proves
\eqref{eq:p213-arbitrary-source-cutoff-fibre-constancy}.

The remaining localization is the one used in the proof of
\cite[Theorem~16.8.5]{CLD25}.  The set $\Omega$ in
\eqref{eq:p213-arbitrary-source-local-omega} is open because the relative
boundary is compact.  The set $U$ misses $\partial(V/W)$, hence is
contained in the relative interior of $V$ and is open in $W$.  It contains
$y$ by \eqref{eq:p213-arbitrary-source-local-boundary}.  Moreover
$F_{\mathrm{trop}}|_U$ is obtained from the proper map
$F_{\mathrm{trop}}:V\to\mathbf R^N$ by base change along
$\Omega\hookrightarrow\mathbf R^N$, and is therefore proper.  Thus the
inverse image of the compact support of $\lambda$ is compact in $U$.
The displayed pullback form consequently extends by zero across
$W\setminus U$, and the equality with $\alpha$ where $\lambda=1$ follows
from \eqref{eq:p213-arbitrary-source-local-form}.

For the last assertion, apply the construction at every point of the
compact test support, starting inside one prescribed cover member that
contains it.  The resulting open neighborhoods cover the support, so
compactness supplies a finite subcover.  Smooth cutoffs subordinate to this
finite cover, with sum one near the compact support, exist by
\cite[Corollary~15.2.9]{CLD25}.
Current identities are local and therefore glue after this finite
decomposition.  If one of those newly chosen cutoffs is retained in a
single local presentation, apply \cite[Corollary~15.1.2]{CLD25} afresh,
append the resulting analytic block $H$, and shrink and refine the finite
cover if necessary.  Compactness again supplies a finite refinement.  At
no point is the cutoff itself used as an analytic coordinate.
\end{proof}
